%==================== RESULTADOS E DISCUSSÃO ====================
\section{Resultados e discussão}

% ESCREVER DEPOIS DA COLETA.
%
% Regra que vale para a seção inteira: nenhum número aqui é digitado
% à mão. As figuras e tabelas são GERADAS pelas rotinas em
% data/analytics, a partir dos dados já anonimizados (P01, P02, ...).
% É o que torna a análise reprodutível e o que permite recompilar o
% documento sem acesso a dado identificável.

\subsection{Caracterização da amostra}

% \input{tables/tab_amostra}
% Gerar com: idade (média, desvio, faixa), curso, período, experiência
% prévia com Realidade Virtual.
[ESCREVER]

\subsection{Medo declarado por estímulo}

% Do questionário inicial (Apêndice B). Gráfico de barras com a média
% do medo declarado (0--10) para cada um dos cinco animais.
% Serve de linha de base para tudo o que vem depois.
[ESCREVER]

\subsection{Progressão pelos níveis}

% Quantos participantes alcançaram cada nível, e quanto tempo
% permaneceram em cada um. O tempo por nível é informativo: parada
% longa pode indicar tanto exploração quanto hesitação.
[ESCREVER]

\subsection{Comportamento durante a exposição}

% As variáveis do Quadro~\ref{tab:variaveis}: aproximação,
% afastamento, agitação, atenção ao estímulo, exploração, interação.
[ESCREVER]

\subsection{Medo declarado frente ao comportamento observado}

% A análise mais interessante do trabalho, e a que melhor responde à
% questão de pesquisa: quem declarou medo alto manteve maior
% distância? Chegou a níveis mais baixos? Olhou menos para o
% estímulo?
%
% É o cruzamento entre o Apêndice B e as medidas automáticas -- entre
% o que a pessoa disse e o que ela fez.
[ESCREVER]

\subsection{Percepção dos participantes}

% Do questionário final (Apêndice C): presença, realismo, facilidade
% de uso, clareza das instruções, desconforto físico, sensação de
% controle sobre o próprio ritmo.
[ESCREVER]

\subsection{Desempenho nos aparelhos utilizados}

% Rodar ferramentas/teste-sistema.html em cada aparelho usado nas
% sessões e registrar aqui. Os números da Seção 5 referem-se a
% renderização por software e não substituem esta medição.
[ESCREVER]

\subsection{Limitações}

% Ser explícita aqui conta a favor na banca. Candidatas já
% identificadas:
%
%   - amostra pequena e por conveniência;
%   - ausência de grupo de controle;
%   - participantes não necessariamente fóbicos -- o trabalho mede
%     percepção e interação, não eficácia;
%   - sessão única, sem acompanhamento ao longo do tempo;
%   - ausência de medida fisiológica;
%   - estímulos pequenos apresentados em escala maior que a real;
%   - sons sintetizados, não gravados.
[ESCREVER]
